%=====================================
%   20260204-rits.tex
%   Microlocal Morse Theory and Truncated Microsupport
%   toshi2019
%=====================================
\documentclass[dvipdfmx,10pt,aspectratio=169,leqno]{beamer}% dvipdfmxしたい
\title{Microlocal Morse Theory and Truncated Microsupport\\
（超局所モース理論と打切り超局所台）
}
%\subtitle{}
\date{2026年2月4日 \\令和7年度数学専攻修士論文審査会}
\author{大柴 寿浩}

\institute[NITech]
{
北海道大学大学院理学院数学専攻\\
指導教員：本多尚文
}


\usetheme{metropolis}
%\usetheme{Singapore}
\usecolortheme{rose}
\usefonttheme{professionalfonts}
\setbeamertemplate{navigaton symbols}{\false}
\setbeamertemplate{footline}[frame number]
%\usepackage{graphicx,xcolor}
\usepackage{appendixnumberbeamer}





\usepackage{bxdpx-beamer}% dvipdfmxなので必要
\usepackage{pxjahyper}% 日本語で'しおり'したい
\usepackage{minijs}% min10ヤダ
\renewcommand{\kanjifamilydefault}{\gtdefault}
\renewcommand{\emph}[1]{{\upshape\bfseries #1}}

\usepackage{amsmath}
\usepackage{amsthm}
\usepackage{tikz}
\usepackage{color}
\usepackage{ascmac}
\usepackage{amsfonts}
\usepackage{mathrsfs}
\usepackage[mathscr]{eucal}
\usepackage{mathtools}
\usepackage{amssymb}
\usepackage{graphicx}
\usepackage{fancybox}
\usepackage{enumerate}
\usepackage{verbatim}
\usepackage{subfigure}
\usepackage{proof}
\usepackage{listings}
\usepackage{otf}
\usepackage[all]{xy}
\usepackage{amscd}
%\usepackage[dvipdfmx]{hyperref}

\usepackage{xcolor}
\definecolor{darkgreen}{rgb}{0,0.45,0} 
\definecolor{darkred}{rgb}{0.75,0,0}
\definecolor{darkblue}{rgb}{0,0,0.6} 
\hypersetup{
    colorlinks=true,
    citecolor=darkgreen,
    linkcolor=darkblue,
    urlcolor=darkblue,
}

\usetikzlibrary{positioning}


\usepackage{latexsym}
\usepackage{wrapfig}
\usepackage{layout}
\usepackage{url}

\usepackage{okumacro}

\usepackage{comment}

%\usepackage{jdash}
%\usepackage{pxjahyper}
% ----------------------------
% commmand
% ----------------------------
% 執筆に便利なコマンド集です. 
% コマンドを追加する場合は下のスペースへ. 

% 集合の記号 (黒板文字)
\newcommand{\NN}{\mathbb{N}}
\newcommand{\ZZ}{\mathbb{Z}}
\newcommand{\QQ}{\mathbb{Q}}
\newcommand{\RR}{\mathbb{R}}
\newcommand{\CC}{\mathbb{C}}
\newcommand{\PP}{\mathbb{P}}
\newcommand{\KK}{\mathbb{K}}


% 集合の記号 (太文字)
\newcommand{\nn}{\mathbf{N}}
\newcommand{\zz}{\mathbf{Z}}
\newcommand{\qq}{\mathbf{Q}}
\newcommand{\rr}{\mathbf{R}}
\newcommand{\cc}{\mathbf{C}}
\newcommand{\pp}{\mathbf{P}}
\newcommand{\kk}{\mathbf{k}}

% 特殊な写像の記号
\newcommand{\ev}{\mathop{\mathrm{ev}}\nolimits} % 値写像
\newcommand{\pr}{\mathop{\mathrm{pr}}\nolimits} % 射影

% スクリプト体にするコマンド
%   例えば　{\mcal C} のように用いる
\newcommand{\mcal}{\mathcal}

% 花文字にするコマンド 
%   例えば　{\h C} のように用いる
\newcommand{\h}{\mathscr}

% ヒルベルト空間などの記号
\newcommand{\F}{\mcal{F}}
\newcommand{\X}{\mcal{X}}
\newcommand{\Y}{\mcal{Y}}
\newcommand{\Hil}{\mcal{H}}
\newcommand{\RKHS}{\Hil_{k}}
\newcommand{\Loss}{\mcal{L}_{D}}
\newcommand{\MLsp}{(\X, \Y, D, \Hil, \Loss)}

% 偏微分作用素の記号
\newcommand{\p}{\partial}

% 角カッコの記号 (内積は下にマクロがあります)
\newcommand{\lan}{\langle}
\newcommand{\ran}{\rangle}



% 圏の記号など
\newcommand{\Set}{{\bf Set}}
\newcommand{\Vect}{{\bf Vect}}
\newcommand{\FDVect}{{\bf FDVect}}
%\newcommand{\Ring}{{\bf Ring}}
\newcommand{\Ab}{{\mathsf{Ab}}}
\newcommand{\Mod}{\mathop{\mathrm{Mod}}\nolimits}
\newcommand{\Modf}{\mathop{\mathrm{Mod}^\mathrm{f}}\nolimits}
\newcommand{\CGA}{{\bf CGA}}
\newcommand{\GVect}{{\bf GVect}}
\newcommand{\Lie}{{\bf Lie}}
\newcommand{\dLie}{{\bf Liec}}



% 射の集合など
\newcommand{\Map}{\mathop{\mathrm{Map}}\nolimits} % 写像の集合
\newcommand{\Hom}{\mathop{\mathrm{Hom}}\nolimits} % 射集合
\newcommand{\End}{\mathop{\mathrm{End}}\nolimits} % 自己準同型の集合
\newcommand{\Aut}{\mathop{\mathrm{Aut}}\nolimits} % 自己同型の集合
\newcommand{\Mor}{\mathop{\mathrm{Mor}}\nolimits} % 射集合
\newcommand{\Ker}{\mathop{\mathrm{Ker}}\nolimits} % 核
\newcommand{\Img}{\mathop{\mathrm{Im}}\nolimits} % 像
\newcommand{\Cok}{\mathop{\mathrm{Coker}}\nolimits} % 余核
\newcommand{\Cim}{\mathop{\mathrm{Coim}}\nolimits} % 余像

% その他便利なコマンド
\newcommand{\dip}{\displaystyle} % 本文中で数式モード
\newcommand{\e}{\varepsilon} % イプシロン
\newcommand{\dl}{\delta} % デルタ
\newcommand{\pphi}{\varphi} % ファイ
\newcommand{\ti}{\tilde} % チルダ
\newcommand{\pal}{\parallel} % 平行
\newcommand{\op}{{\rm op}} % 双対を取る記号
\newcommand{\lcm}{\mathop{\mathrm{lcm}}\nolimits} % 最小公倍数の記号
\newcommand{\Probsp}{(\Omega, \F, \P)} 
\newcommand{\argmax}{\mathop{\rm arg~max}\limits}
\newcommand{\argmin}{\mathop{\rm arg~min}\limits}





% ================================
% コマンドを追加する場合のスペース 
%\newcommand{\OO}{\mcal{O}}



\makeatletter
\renewenvironment{proof}[1][\proofname]{\par
  \pushQED{\qed}%
  \normalfont \topsep6\p@\@plus6\p@\relax
  \trivlist
  \item[\hskip\labelsep
%        \itshape
         \bfseries
%    #1\@addpunct{.}]\ignorespaces
    {#1}]\ignorespaces
}{%
  \popQED\endtrivlist\@endpefalse
}
\makeatother

\renewcommand{\proofname}{\textrm{Proof.}}



%\renewcommand\proofname{\bf 証明} % 証明
\numberwithin{equation}{subsection}
\newcommand{\cTop}{\textsf{Top}}
%\newcommand{\cOpen}{\textsf{Open}}
\newcommand{\Op}{\mathop{\textsf{Op}}\nolimits}
\newcommand{\Ob}{\mathop{\textrm{Ob}}\nolimits}
\newcommand{\id}{\mathop{\mathrm{id}}\nolimits}
\newcommand{\pt}{\mathop{\mathrm{pt}}\nolimits}
\newcommand{\res}{\mathop{\rho}\nolimits}
\newcommand{\A}{\mcal{A}}
\newcommand{\B}{\mcal{B}}
\newcommand{\C}{\mcal{C}}
\newcommand{\D}{\mcal{D}}
\newcommand{\E}{\mcal{E}}
\newcommand{\G}{\mcal{G}}
%\newcommand{\H}{\mcal{H}}
\newcommand{\I}{\mcal{I}}
\newcommand{\J}{\mcal{J}}
\newcommand{\OO}{\mcal{O}}
\newcommand{\Ring}{\mathop{\textsf{Ring}}\nolimits}
\newcommand{\cAb}{\mathop{\textsf{Ab}}\nolimits}
%\newcommand{\Ker}{\mathop{\mathrm{Ker}}\nolimits}
\newcommand{\im}{\mathop{\mathrm{Im}}\nolimits}
\newcommand{\Coker}{\mathop{\mathrm{Coker}}\nolimits}
\newcommand{\Coim}{\mathop{\mathrm{Coim}}\nolimits}
\newcommand{\rank}{\mathop{\mathrm{rank}}\nolimits}
\newcommand{\Ht}{\mathop{\mathrm{Ht}}\nolimits}
\newcommand{\supp}{\mathop{\mathrm{supp}}\nolimits}
\newcommand{\colim}{\mathop{\mathrm{colim}}}
\newcommand{\Tor}{\mathop{\mathrm{Tor}}\nolimits}

\newcommand{\cat}{\mathscr{C}}

%筆記体
\newcommand{\cA}{\mcal{A}}
\newcommand{\cB}{\mcal{B}}
\newcommand{\cC}{\mcal{C}}
\newcommand{\cD}{\mcal{D}}
\newcommand{\cE}{\mcal{E}}
\newcommand{\cF}{\mcal{F}}
\newcommand{\cG}{\mcal{G}}
\newcommand{\cH}{\mcal{H}}
\newcommand{\cI}{\mcal{I}}
\newcommand{\cJ}{\mcal{J}}
\newcommand{\cK}{\mcal{K}}
\newcommand{\cL}{\mcal{L}}
\newcommand{\cM}{\mcal{M}}
\newcommand{\cN}{\mcal{N}}
\newcommand{\cO}{\mcal{O}}
\newcommand{\cP}{\mcal{P}}
\newcommand{\cQ}{\mcal{Q}}
\newcommand{\cR}{\mcal{R}}
\newcommand{\cS}{\mcal{S}}
\newcommand{\cT}{\mcal{T}}
\newcommand{\cU}{\mcal{U}}
\newcommand{\cV}{\mcal{V}}
\newcommand{\cW}{\mcal{W}}
\newcommand{\cX}{\mcal{X}}
\newcommand{\cY}{\mcal{Y}}
\newcommand{\cZ}{\mcal{Z}}


\newcommand{\scA}{\mathscr{A}}
\newcommand{\scB}{\mathscr{B}}
\newcommand{\scC}{\mathscr{C}}
\newcommand{\scD}{\mathscr{D}}
\newcommand{\scE}{\mathscr{E}}
\newcommand{\scF}{\mathscr{F}}
\newcommand{\scN}{\mathscr{N}}
\newcommand{\scO}{\mathscr{O}}
\newcommand{\scV}{\mathscr{V}}
\newcommand{\scU}{\mathscr{U}}


\newcommand{\ibA}{\mathop{\text{\textit{\textbf{A}}}}}
\newcommand{\ibB}{\mathop{\text{\textit{\textbf{B}}}}}
\newcommand{\ibC}{\mathop{\text{\textit{\textbf{C}}}}}
\newcommand{\ibD}{\mathop{\text{\textit{\textbf{D}}}}}
\newcommand{\ibE}{\mathop{\text{\textit{\textbf{E}}}}}
\newcommand{\ibF}{\mathop{\text{\textit{\textbf{F}}}}}
\newcommand{\ibG}{\mathop{\text{\textit{\textbf{G}}}}}
\newcommand{\ibH}{\mathop{\text{\textit{\textbf{H}}}}}
\newcommand{\ibI}{\mathop{\text{\textit{\textbf{I}}}}}
\newcommand{\ibJ}{\mathop{\text{\textit{\textbf{J}}}}}
\newcommand{\ibK}{\mathop{\text{\textit{\textbf{K}}}}}
\newcommand{\ibL}{\mathop{\text{\textit{\textbf{L}}}}}
\newcommand{\ibM}{\mathop{\text{\textit{\textbf{M}}}}}
\newcommand{\ibN}{\mathop{\text{\textit{\textbf{N}}}}}
\newcommand{\ibO}{\mathop{\text{\textit{\textbf{O}}}}}
\newcommand{\ibP}{\mathop{\text{\textit{\textbf{P}}}}}
\newcommand{\ibQ}{\mathop{\text{\textit{\textbf{Q}}}}}
\newcommand{\ibR}{\mathop{\text{\textit{\textbf{R}}}}}
\newcommand{\ibS}{\mathop{\text{\textit{\textbf{S}}}}}
\newcommand{\ibT}{\mathop{\text{\textit{\textbf{T}}}}}
\newcommand{\ibU}{\mathop{\text{\textit{\textbf{U}}}}}
\newcommand{\ibV}{\mathop{\text{\textit{\textbf{V}}}}}
\newcommand{\ibW}{\mathop{\text{\textit{\textbf{W}}}}}
\newcommand{\ibX}{\mathop{\text{\textit{\textbf{X}}}}}
\newcommand{\ibY}{\mathop{\text{\textit{\textbf{Y}}}}}
\newcommand{\ibZ}{\mathop{\text{\textit{\textbf{Z}}}}}

\newcommand{\ibx}{\mathop{\text{\textit{\textbf{x}}}}}

%\newcommand{\Comp}{\mathop{\mathrm{C}}\nolimits}
%\newcommand{\Komp}{\mathop{\mathrm{K}}\nolimits}
%\newcommand{\Domp}{\mathop{\mathsf{D}}\nolimits}%複体のホモトピー圏
%\newcommand{\Comp}{\mathrm{C}}
%\newcommand{\Komp}{\mathrm{K}}
%\newcommand{\Domp}{\mathsf{D}}%複体のホモトピー圏

\newcommand{\Comp}{\mathsf{C}}
\newcommand{\Komp}{\mathsf{K}}
\newcommand{\Domp}{\mathsf{D}}
\newcommand{\Kompl}{\mathop{\mathsf{K}^\mathrm{+}}\nolimits}
\newcommand{\Kompu}{\mathop{\mathsf{K}^\mathrm{-}}\nolimits}
\newcommand{\Kompb}{\mathop{\mathsf{K}^\mathrm{b}}\nolimits}
\newcommand{\Dompl}{\mathop{\mathsf{D}^\mathrm{+}}\nolimits}
\newcommand{\Dompu}{\mathop{\mathsf{D}^\mathrm{-}}\nolimits}
\newcommand{\Dompb}{\mathop{\mathsf{D}^\mathrm{b}}\nolimits}
\newcommand{\Dompbf}{\mathop{\mathsf{D}_\mathrm{f}^\mathrm{b}}\nolimits}




\newcommand{\CCat}{\Comp(\cat)}
\newcommand{\KCat}{\Komp(\cat)}
\newcommand{\DCat}{\Domp(\cat)}%圏Cの複体のホモトピー圏
\newcommand{\HOM}{\mathop{\mathscr{H}\hspace{-2pt}om}\nolimits}%内部Hom
\newcommand{\RHOM}{\mathop{\mathrm{R}\hspace{-1.5pt}\HOM}\nolimits}

%\newcommand{\muS}{\mathop{\mathrm{SS}}\nolimits}
\newcommand{\muS}[1][]{\mathop{\mathrm{SS}}\nolimits_{#1}}
\newcommand{\RG}{\mathop{\mathrm{R}\hspace{-0pt}\Gamma}\nolimits}
\newcommand{\RHom}{\mathop{\mathrm{R}\hspace{-1.5pt}\Hom}\nolimits}
%\newcommand{\Rder}{\mathrm{R}}

\newcommand{\Rder}[1][]{\mathop{\mathrm{R}^{#1}\hspace{-2pt}}\nolimits}


\newcommand{\simar}{\mathrel{\overset{\sim}{\rightarrow}}}%同型右矢印
\newcommand{\simarr}{\mathrel{\overset{\sim}{\longrightarrow}}}%同型右矢印
\newcommand{\simra}{\mathrel{\overset{\sim}{\leftarrow}}}%同型左矢印
\newcommand{\simrra}{\mathrel{\overset{\sim}{\longleftarrow}}}%同型左矢印

\newcommand{\hocolim}{{\mathrm{hocolim}}}
\newcommand{\indlim}[1][]{\mathop{\varinjlim}\limits_{#1}}
\newcommand{\sindlim}[1][]{\smash{\mathop{\varinjlim}\limits_{#1}}\,}
\newcommand{\Pro}{\mathrm{Pro}}
\newcommand{\Ind}{\mathrm{Ind}}
\newcommand{\prolim}[1][]{\mathop{\varprojlim}\limits_{#1}}
\newcommand{\sprolim}[1][]{\smash{\mathop{\varprojlim}\limits_{#1}}\,}

\newcommand{\Sh}{\mathrm{Sh}}
\newcommand{\PSh}{\mathrm{PSh}}

\newcommand{\rmD}{\mathrm{D}}

\newcommand{\Lloc}[1][]{\mathord{\mathcal{L}^1_{\mathrm{loc},{#1}}}}
\newcommand{\ori}{\mathord{\mathrm{or}}}
\newcommand{\Db}{\mathord{\mathscr{D}b}}

\newcommand{\codim}{\mathop{\mathrm{codim}}\nolimits}



\newcommand{\gld}{\mathop{\mathrm{gld}}\nolimits}
\newcommand{\wgld}{\mathop{\mathrm{wgld}}\nolimits}


\newcommand{\tens}[1][]{\mathbin{\otimes_{\raise1.5ex\hbox to-.1em{}{#1}}}}
\newcommand{\etens}{\mathbin{\boxtimes}}
\newcommand{\ltens}[1][]{\mathbin{\overset{\mathrm{L}}\tens}_{#1}}
\newcommand{\mtens}[1][]{\mathbin{\overset{\mathrm{\mu}}\tens}_{#1}}
\newcommand{\lltens}[1][]{{\mathop{\tens}\limits^{\mathrm{L}}_{#1}}}
\newcommand{\letens}{\overset{\mathrm{L}}{\etens}}
\newcommand{\detens}{\underline{\etens}}
\newcommand{\ldetens}{\overset{\mathrm{L}}{\underline{\etens}}}
\newcommand{\dtens}[1][]{{\overset{\mathrm{L}}{\underline{\otimes}}}_{#1}}

\newcommand{\blk}{\mathord{\ \cdot\ }}
\newcommand{\mres}[2][]{{\left.{#1}\right\rvert}_{#2}}


%\newcommand{\hocolim}{{\mathrm{hocolim}}}
%\newcommand{\indlim}[1][]{\mathop{\varinjlim}\limits_{#1}}
%\newcommand{\sindlim}[1][]{\smash{\mathop{\varinjlim}\limits_{#1}}\,}
%\newcommand{\Pro}{\mathrm{Pro}}
%\newcommand{\Ind}{\mathrm{Ind}}
%\newcommand{\prolim}[1][]{\mathop{\varprojlim}\limits_{#1}}
%\newcommand{\sprolim}[1][]{\smash{\mathop{\varprojlim}\limits_{#1}}\,}
\newcommand{\proolim}[1][]{\mathop{\text{\rm``{$\varprojlim$}''}}\limits_{#1}}
\newcommand{\sproolim}[1][]{\smash{\mathop{\rm``{\varprojlim}''}\limits_{#1}}}
\newcommand{\inddlim}[1][]{\mathop{\text{\rm``{$\varinjlim$}''}}\limits_{#1}}
\newcommand{\sinddlim}[1][]{\smash{\mathop{\text{\rm``{$\varinjlim$}''}}\limits_{#1}}\,}
\newcommand{\ooplus}{\mathop{\text{\rm``{$\oplus$}''}}\limits}
\newcommand{\bbigsqcup}{\mathop{``\bigsqcup''}\limits}
\newcommand{\bsqcup}{\mathop{``\sqcup''}\limits}
\newcommand{\dsum}[1][]{\mathbin{\oplus_{#1}}}

\newcommand{\Fct}{\mathop{\mathsf{Fct}}\nolimits}
\newcommand{\Red}{\mathop{\mathsf{Red}}\nolimits}
\newcommand{\Cone}{\mathop{\mathsf{Cone}}\nolimits}
\newcommand{\epi}{\mathop{\mathsf{epi}}\nolimits}




\newcommand{\mtan}[1][]{T\hspace{-1.75pt}{#1}}
\newcommand{\mtp}[1][]{{}^{t\!}{#1}} % transpose of the morphism
%\newcommand{\mpb}[1][]{{}^{t\!}{#1}'} % pull-back of the morphism
\newcommand{\mpb}[1][]{{#1}_{\pi}} % pull-back of the morphism
\newcommand{\mdiff}[1][]{{#1}_{d}} % pull-back of the morphism
\newcommand{\mptan}[1][]{{#1}_{\tau}} % pull-back of the morphism


%\newcommand{\pb}[1][]{\mathbin{\mathop{\times}\limits_{#1}}}

\newcommand{\cotb}[1][]{T^{\ast}{#1}}
\newcommand{\conm}[2][]{T_{#1}^{\hspace{0.7pt}\ast}{#2}}
\newcommand{\norb}[2][]{T_{#1}{#2}}

\newcommand{\cotz}[1][]{T^{\ast}_{#1}{#1}}

\newcommand{\cotn}[1][]{\dot{T}^{\ast}{#1}}


\newcommand{\muhom}{\mu hom}
\newcommand{\muHom}[1][]{\mathrm{Hom}^\mu_{\raise1.5ex\hbox to.1em{}#1}}

\newcommand{\mucat}[2][]{\mathsf{D}^{\mathrm{b}}_{#1}(#2)}
%\newcommand{\mucat}[2][]{\mathsf{D}^{\mathrm{b}}_{#1}(#2)}

\newcommand{\conv}[1][]{\mathbin{\mathop{\circ}\limits_{#1}}}

\newcommand{\pb}[1][]{\mathbin{\mathop{\times}\limits_{#1}}}

\newcommand{\Sol}[2][]{\mathop{\mathcal{S}ol}\nolimits_{#1}(#2)}



%================================================
% 自前の定理環境
%   https://mathlandscape.com/latex-amsthm/
% を参考にした
\newtheoremstyle{mystyle}%   % スタイル名
    {5pt}%                   % 上部スペース
    {5pt}%                   % 下部スペース
    {}%              % 本文フォント
    {}%                  % 1行目のインデント量
    {\bfseries}%                      % 見出しフォント
    {.}%                     % 見出し後の句読点
    {12pt}%                     % 見出し後のスペース
    {\thmname{#1}\thmnumber{ #2}\thmnote{{\hspace{2pt}\normalfont (#3)}}}% % 見出しの書式

\theoremstyle{mystyle}
\newtheorem{AXM}{公理}%[section]
\newtheorem{DFN}[AXM]{定義}
\newtheorem{THM}[AXM]{定理}
\newtheorem{HM}[AXM]{定理もどき}
\newtheorem*{THM*}{定理}
\newtheorem{PRP}[AXM]{命題}
\newtheorem{LMM}[AXM]{補題}
\newtheorem{CRL}[AXM]{系}
\newtheorem{EG}[AXM]{例}
\newtheorem*{EG*}{例}
\newtheorem{RMK}[AXM]{注意}
\newtheorem{CNV}[AXM]{約束}
\newtheorem{CNJ}[AXM]{予想}
\newtheorem{CMT}[AXM]{コメント}
\newtheorem*{CMT*}{コメント}
\newtheorem{NTN}[AXM]{記号}
\newtheorem{CLM}[AXM]{Claim}
\newtheorem{Prop}[AXM]{Proposition}

% 定理環境ここまで
%====================================================

\usepackage{framed}
\definecolor{lightgray}{rgb}{0.75,0.75,0.75}
\renewenvironment{leftbar}{%
  \def\FrameCommand{\textcolor{lightgray}{\vrule width 4pt} \hspace{10pt}}% 
  \MakeFramed {\advance\hsize-\width \FrameRestore}}%
{\endMakeFramed}
\newenvironment{redleftbar}{%
  \def\FrameCommand{\textcolor{lightgray}{\vrule width 1pt} \hspace{10pt}}% 
  \MakeFramed {\advance\hsize-\width \FrameRestore}}%
 {\endMakeFramed}



\renewcommand{\Re}{\mathop{\mathrm{Re}}\nolimits}


% =================================





% ---------------------------
% new definition macro
% ---------------------------
% 便利なマクロ集です

% 内積のマクロ
%   例えば \inner<\pphi | \psi> のように用いる
\def\inner<#1>{\langle #1 \rangle}

% ================================
% マクロを追加する場合のスペース 

%=================================







\def\fracinline#1/#2{\mbox{\raise0.5ex\hbox{\footnotesize$#1$}{\hskip-.1em$/$\hskip-.1em}\raise-0.5ex\hbox{\footnotesize$#2$}}}

\begin{document}




\begin{frame}
    \titlepage
    %\insertsubtitle
\end{frame}
\begin{comment}
    \begin{frame}\frametitle{書くこと}
    \begin{itemize}
        \item 特に興味を持った定理や理論の背景
        \item それを記述するための記号や概念の導入
        \item 正確な主張の紹介
        \item 証明の基本的なアイデアやアウトラインの紹介
        \item または定理の過程を満たす具体例や定理の応用例の紹介
    \end{itemize}
    \end{frame}
\end{comment}
%\setcounter{section}{3}
%\section*{はじめに}
\begin{frame}
    \frametitle{研究の概要}
    \begin{center}
      動機：「層のコホモロジーの延長可能性を調べたい」
    \end{center}
    \vspace{10pt}
    \[\vcenter{\xymatrix@C=42pt@R=26pt{
        \txt{\(\muS[](F)\)\\
        \cite{KS82}}
        %\ar[dd]
        &
        \txt{超局所モースの補題\\
        \cite{KS90}}
        %\ar[dd]
        &
        \txt{非特性変形補題\\
        \cite{KS90}}
        %\ar[dd]
        \\
        %\ar[r] 
        &
        &
        %\ar[l]
        \\
        \txt{\(\muS[k](F)\)\\
        \cite{KFS03a}}
        &
        \txt{打切り超局所モースの補題\\主結果1}
        &\txt{打切り非特性変形補題\\主結果2}
    }}\]
\end{frame}
%\begin{frame}
%    \frametitle{発表内容}
%    \tableofcontents
%\end{frame}

%\section[背景]{背景}
%\subsection{層理論}
\begin{comment}
\begin{frame}
    \frametitle{（前）層の定義}
    \begin{DFN}[層]\label{def:sheaf}
        位相空間\(X\)上の\(\kk\)加群の前層\(F\)とは，
        \(\kk\)加群の族と\(\kk\)加群の射の族の組\[\left(
        \left(F(U)\right)_{U\subset X},
        \left(\rho_{VU}\colon F(U)\to F(V)\right)_{(V\hookrightarrow U)}
        \right)\]で次をみたすもののこと．
        簡単のため\(\kk\)は\(\zz\)か体とする．
        \vspace*{-5pt}
        \begin{enumerate}
            %\renewcommand{\labelenumi}{({\arabic{enumi}})}
            \item $\rho_{UU} = \id_{F(U)}$. 
            \item $W \subset V \subset U$ ならば, 
            $\rho_{WU} = \rho_{WV} \circ \rho_{VU}$.\label{enum:res}
        \end{enumerate}
        \vspace*{-5pt}
        \(F\)が\textbf{層}であるとは，
        \(X\)の任意の開集合\(U\)と
        その開被覆\(\left(U_i\right)_{i\in I}\)に対して，
        次が成り立つこと．
        \vspace*{-5pt}
        \begin{enumerate}[(S1)]
            \setcounter{enumi}{-1}
            \setlength{\leftskip}{10pt}
            \item \(F(\varnothing)=0\).
            \item \(s\in F(U)\)が各\(i\in I\)に対して\(U_i\)上\(s\rvert_{U_{i}}=0\)ならば，
            \(s=0\).\label{sheaf-cond1}
            \item \(\left(s_i\right)_i\in\prod_{i}F(U_i)\)が
            各\(i,j\in I\)で\(U_i\cap U_j\ne\varnothing\)となるものに対して\(U_i\cap U_j\)上\(s_i\rvert_{U_i\cap U_j}-s_j\rvert_{U_i\cap U_j}=0\)ならば，
            \(s\in F(U)\)で各\(U_i\)上\(s\rvert_i=s_i\)となるものが存在する．\label{sheaf-cond2}
        \end{enumerate}
    \end{DFN}
\end{frame}

\begin{frame}
    \frametitle{層の導来圏}
    導来圏の理論を用いると見通しがよい
    \begin{NTN}
        層の有界な複体を対象とし，
        コホモロジー間の同型を誘導する射（擬同型）を可逆にした
        圏を\(\Dompb(\kk_X)\)で表す．
    \end{NTN}
    \(\Dompb(\cc_X)\)では例えば定数層のド・ラーム分解が同型になる：
    \[\vcenter{\xymatrix{
        0\ar[r]&
        \cc_X \ar[r] \ar[d]^-{\mathrm{qis}}&
        0\\
        0\ar[r]&
        \Omega_X^0\ar[r]&
        \Omega_X^1\ar[r]&
        \cdots\ar[r]&
        \Omega_X^n\ar[r]&0
    }}\]
\end{frame}
\end{comment}

%\subsection{超局所層理論}
\newcommand{\Char}{\mathop{\mathrm{Ch}}\nolimits}
\newcommand{\singsupp}{\mathop{\mathrm{singsupp}}\nolimits}
\newcommand{\SingSpec}{\mathop{\mathrm{S.S.}}\nolimits}
%\newcommand{\SingSpec}{\mathop{\mathrm{WF}_{\mathrm{A}}}\nolimits}

\begin{comment}
\begin{frame}\frametitle{``超局所''？--- 超局所解析からの復習}
    \[
        \text{``超局所''}=\text{余接束上で局所}
    \]
    いったん，古典的な微分方程式の問題で説明する．
    \begin{itemize}
        \item 余接束：開集合\(\varOmega\subset \rr^n\)に対し
        \(\dot{\pi}\colon\cotn[\varOmega]\coloneqq \cotb[\varOmega]-\cotz[\varOmega]\cong \varOmega\times(\rr^n-\{0\})\to \varOmega\). 
        \item 超関数\(u\)に対し特異台\(
            \singsupp(u)\coloneqq\left\{
                x\in\varOmega;\ \text{\(u(x)\)は\(x\)で解析的でない}
            \right\}\). 
        \item 微分作用素\(P\)に対し\(
            \Char(P)\coloneqq\{(x;\xi)\in\cotn[\varOmega];\ 
                \sigma(P)(x;\xi)=0\}.
            \)\\（ただし\(\sigma(P)\)は\(P\)の``主要表象''）
        \item \(u\)の特異性スペクトル：\(\SingSpec(u)\subset\cotn[\varOmega]\)は\(
            \pi(\SingSpec(u))=\singsupp(u)
        \)をみたす集合．
    \end{itemize}
    それぞれ
    \(\singsupp{u}\subset\varOmega\)
    と
    \(\Char(P), \SingSpec(u)\subset \cotn[\varOmega]\)に注意．
\end{frame}

\begin{frame}\frametitle{超局所解析における定理} 
    \begin{THM}[佐藤の基本定理{\cite{Sato69,Sato70}}]
        \(f\)：既知関数，
        \(P\)：実解析的係数微分作用素とする．
        \(u\)が微分方程式\(Pu=f\)の超関数解のとき
        \[
            \SingSpec(u)\subset \Char(P)\cup\SingSpec(f).
        \]
    \end{THM}
    これを\(f=0\), \(P\)：楕円型の場合に適用すると
    \[
        \Char(P)\cup \SingSpec(f)=\varnothing
    \]
    なので\(\singsupp(u)=\pi(\SingSpec(u))=\varnothing\)．
    すなわち，
    \begin{alertblock}{楕円型正則性定理}
        実解析的係数の楕円型方程式の解は解析的である．
    \end{alertblock}
    このように，底空間の対象を余接束上に持ち上げる（超局所的に考える）と，
    その対象の特異性が分析できる．
\end{frame}
\end{comment}



\begin{frame}
    \frametitle{\(\text{「超局所」}=\text{「余接束上で局所」}\)}
    \begin{exampleblock}{超局所解析における例}
        微分方程式\(Pu=0\)を考える．余接束の部分集合
        \begin{itemize}
            \item 超関数\(u\)の波面集合（特異性スペクトル）
            \item 線形微分作用素\(P\)の特性多様体
        \end{itemize}
        を用いると，解\(u\)の特異性が分析できる（楕円型正則性など）．
    \end{exampleblock}
    以上の層理論での類似が超局所層理論．
    層には次のような例がある．
    \begin{itemize}
      \item \(C^{\infty}\)関数の層\(\cC^{\infty}_{M}\)（\(M\)：実\(C^{\infty}\)多様体）
      \item 微分方程式系\(\cM\)の正則関数解の層\(\Sol[X]{\cM}\coloneqq\RHOM_{\cD_{X}}(\cM,\cO_{X})\)（\(X\)：複素多様体）
    \end{itemize}
    多様体\(M\)上の層（の複体）\(F\in\Dompb(\kk_{M})\)に対し，
    余接束の部分集合\(\muS[](F)\subset \cotb[M]\)を定める．
\end{frame}

\begin{frame}\frametitle{層の超局所台 --- 延長可能性の判定}
    \begin{DFN}[層の超局所台]
        \(p\notin \muS(F)\)であるとは，
        \(p\)の\(\cotb[M]\)における開近傍\(U\)で，任意の\(x_1\in X\)と
        \(x_1\)の近傍上で定義された\(C^{1}\)関数\(\varphi\)に対し，
        次をみたすものが存在することをいう．
        \[
            \varphi(x_1)=0,\ d\varphi(x_1)\in U
            \Longrightarrow
            H^{j}_{\{\varphi\geqq0\}}(F)_{x_1}=0
            \ (j\in\zz).
        \]
    \end{DFN}
    微分方程式との関係は次の通り．ただし\(\cM=\cD_{X}/\cD_{X}P\)とおいた．
    \[\vcenter{\xymatrix@C=36pt@R=20pt{
      \txt{\(Pu=0\)の解}
      \ar@{.>}[r]
      \ar@{<->}[d]
      &
      \txt{\(x\)で\(\xi\)方向に一意に延長可能}
      &
      \ldots\txt{微分方程式}\\
      \Sol[X]{\cM}
      \ar@{|->}[r]
      &
      \muS[](\Sol[X]{\cM})\not\ni (x;\xi)
      \ar@{<=>}[u]
      &
      \ldots\txt{超局所層理論}
    }}\]
    %\(\Dompb(\kk_X)\ni F\mapsto \SS(F)\subset T^{\ast}X\)    
    %超局所台を用いると，関数の臨界点を層係数に一般化することができる．
    %\begin{DFN}
    %    \begin{itemize}
    %        \item \(x\in M\)が\(\varphi\colon M\to \rr\)の臨界点
    %        であるとは\(d\varphi(x)=0\)となること．
    %        \item \(x\in M\)が\(\varphi\colon M\to \rr\)の
    %        \textcolor{brown}{\(F\)に関する}臨界点であるとは\(d\varphi(x)\in\muS[](F)\)となること．
    %    \end{itemize}
    %\end{DFN}
    \begin{alertblock}{注意}
        超局所台は，\textcolor{red}{すべての次数}のコホモロジーが延長できるか否かを判定する概念．
    \end{alertblock}
    %\vspace{60pt}
\end{frame}

\begin{frame}\frametitle{打切り超局所台 --- \(k\)次までの特異性}
    \begin{alertblock}{注意}
        超局所台は，\textcolor{red}{すべての次数}のコホモロジーが延長できるか否かを判定する概念．
    \end{alertblock}
    \begin{alertblock}{問題点}
      コホモロジーの次数に依存する伝播現象は超局所台\(\muS[](F)\)では
      捉えられない．
    \end{alertblock}
    
    そこで，超局所台を精密化することで次の定義を得る．
    \begin{DFN}[{\cite[Definition 3.4]{KFS03a}}]
        \textcolor{brown}{\(k\in\zz\)とする．}
        \(p\notin \muS[\textcolor{brown}{k}](F)\)であるとは，
        \(p\)の\(\cotb[M]\)における錐状開近傍\(U\)で，
        任意の\(x_1\in \pi(U)\)と
        \(x_1\)の近傍上で定義された\(C^{1}\)関数\(\varphi\)に対し，
        次をみたすものが存在することをいう．
        \[
            \varphi(x_1)=0,\ d\varphi(x_1)\in U
            \Longrightarrow
            H^{j}_{\{\varphi\geqq0\}}(F)_{x_1}=0
            \quad
            (\textcolor{brown}{j\leqq k}).
        \]
    \end{DFN}
\end{frame}



%\section{主結果}
\begin{frame}
    \frametitle{主結果1の元の主張}
    設定は次のとおり．
    \begin{itemize}
        \item \(M\): \(C^{\infty}\)多様体．
        \item \(F\in\Dompb(\kk_{M})\)（\(M\)上の層の複体）.
        \item \(\psi\colon M\to \rr\)：\(C^1\)級関数で
        \(\supp(F)\)上固有なものと\(t\in\rr\)に対し，\(M_{t}\coloneqq \psi^{-1}((-\infty,t))\).
    \end{itemize}
    \begin{THM}[超局所モースの補題 {\cite{KS90}}]
        \(a\leqq\psi(x)< b\)ならば
        \(d\psi(x)\notin\muS[](F)\)であるとする．
        このとき制限射\[
            H^{j}(M_b;F)\longrightarrow H^{j}(M_a;F)
        \]はすべての\(j\in\zz\)について同型である．
    \end{THM}    
    これを\(\muS[k](F)\)を用いて拡張したのが次の主結果1.
\end{frame}

\begin{frame}
    \frametitle{主結果1}
    設定は先ほどと同じ．
    \begin{itemize}
        \item \(M\): \(C^{\infty}\)多様体．
        \item \(F\in\Dompb(\kk_{M})\)（\(M\)上の層の複体）.
        \item \(\psi\colon M\to \rr\)：\(C^1\)級関数で
        \(\supp(F)\)上固有なものと\(t\in\rr\)に対し，\(M_{t}\coloneqq \psi^{-1}((-\infty,t))\).
    \end{itemize}
    \begin{THM}[大柴, 打切り版超局所モースの補題]
        \textcolor{brown}{\(k\in\zz\)とする．}
        \(a\leqq\psi(x)< b\)ならば
        \(d\psi(x)\notin\muS[\textcolor{brown}{k}](F)\)であるとする．
        このとき制限射\[
            H^{j}(M_b;F)\longrightarrow H^{j}(M_a;F)
        \]は\textcolor{brown}{\(j<k\)ならば同型}であり，
        \textcolor{brown}{\(j=k\)ならば単射}である．
    \end{THM}    
    \(k\gg0\)とすれば，元の主張が復元される．

    主結果1はもっと一般の場合に成り立つ，もう一つの主結果から従う．
\end{frame}
\begin{frame}
    \frametitle{主結果2の元の結果}
    %主結果2は次のオリジナルの結果を打切り版に精密化したもの．
  \begin{THM}[非特性変形補題 {\cite{KS90}}]
    \(X\)をハウスドルフ空間とする．
    \(F\in\Dompl(\kk_X)\)とし，
    \(\left(U_{t}\right)_{t\in\rr}\)を\(X\)の開集合の族とする．
    %\(k\in\zz\)とする．
    \(
      Z_{s}\coloneqq \bigcap_{t>s}\overline{U_t\setminus U_s}\ (s\in\rr)
    \)とおく．
    次を仮定する．%\vspace{-8pt}
    \begin{enumerate}[(i)]
      \setlength{\leftskip}{10pt}
      \item \label{item:nonchar1}
      任意の\(t\in\rr\)に対し\(U_{t}=\bigcup_{s<t}U_{s}\)である．
      \item \label{item:nonchar2}
      任意の実数\(s\leqq t\)に対し\(
        \overline{U_{t}\setminus U_{s}}\cap \supp{F}
      \)はコンパクトである．
      \item \label{item:nonchar3}
      任意の実数\(s\leqq t\)と任意の点\(x\in Z_{s}\setminus U_{t}\)に対し，
      \(
        (H^{j}_{X\setminus U_{t}}(F))_{x}=0\quad (j\in\zz). 
      \)
    \end{enumerate}
    このとき，任意の\(t\in\rr\)とすべての\(j\in\zz\)に対し，
    次の射は同型である．\[
      H^{j}\left(\bigcup_{s\in\rr}U_{s};F\right)
      \to
      H^{j}\left(U_{t};F\right).
    \]
  \end{THM}
\end{frame}

\begin{frame}
    \frametitle{期待される打切り版}
  \begin{HM}[打切り版？非特性変形補題]
    \(X\)をハウスドルフ空間とする．
    \(F\in\Dompl(\kk_X)\)とし，
    \(\left(U_{t}\right)_{t\in\rr}\)を\(X\)の開集合の族とする．
    \textcolor{brown}{\(k\in\zz\)とする．}
    \(
      Z_{s}\coloneqq \bigcap_{t>s}\overline{U_t\setminus U_s}\ (s\in\rr)
    \)とおく．
    次を仮定する．%\vspace{-8pt}
    \begin{enumerate}[(i)]
      \setlength{\leftskip}{10pt}
      \item \label{item:t-nonchar1}
      任意の\(t\in\rr\)に対し\(U_{t}=\bigcup_{s<t}U_{s}\)である．
      \item \label{item:t-nonchar2}
      任意の実数\(s\leqq t\)に対し\(
        \overline{U_{t}\setminus U_{s}}\cap \supp{F}
      \)はコンパクトである．
      \item \label{item:t-nonchar3}
      任意の実数\(s\leqq t\)と
      任意の点\(x\in Z_{s}\setminus U_{t}\)に対し，
      \textcolor{brown}{\(j\leqq k\)ならば}
      \(
        (H^{j}_{X\setminus U_{t}}(F))_{x}=0. 
      \)
    \end{enumerate}
    このとき，任意の\(t\in\rr\)に対し，
    次の射は\textcolor{brown}{\(j<k\)ならば同型}であり，
    \textcolor{brown}{\(j=k\)ならば単射}である．
    \[
      H^{j}\left(\bigcup_{s\in\rr}U_{s};F\right)
      \to
      H^{j}\left(U_{t};F\right).
    \]
  \end{HM}
\end{frame}

\begin{frame}
  \frametitle{うまくいかない例}
  \begin{alertblock}{問題点}
    条件\eqref{item:t-nonchar1}--\eqref{item:t-nonchar3}を
    みたしているのにうまくいかないことがある．
  \end{alertblock}
    \(X=\rr^2\)とする．実数\(t\)に対し関数\(\varphi_{t}\)を次で定める．
    \[
      \varphi_{t}(x)\coloneqq\begin{cases}
        \max\left(0,t(2-\left|x\right|)\right)&(\left|x\right| < 2),\\
        0& (\left|x\right|\geqq 2).
      \end{cases}
    \]
    このとき\(X\)の開集合\(U_{t}\)を
    \(
      U_{t}\coloneqq
      \left\{
        (x,y)\in X;\ y<\varphi_{t}(x)
      \right\}
    \)とする．次ページのように層\(\kk_{Z}\)を定めると，
    \((U_{t})_{t}\)と\(\kk_{Z}\)は\eqref{item:t-nonchar1}--\eqref{item:t-nonchar3}をみたす．
      \begin{figure}[htb]
      \centering
      \begin{tikzpicture}[scale=0.50]
        \def\t{0.7}
        \def\u{1.4}
        \def\xmax{3.5}
        \def\ymin{-1}
        \def\ymax{3.3}

        %------------------
        % 塗りつぶし
        %------------------

        % y<0 の部分
        \fill[green!30] (-\xmax,0) rectangle (\xmax,\ymin);
        \fill[blue!25] (-\xmax,0) rectangle (\xmax,\ymin);

        % 三角形部分 y < t(2-|x|)
        \fill[green!30]
          (-2,0) -- (0,{2*\u}) -- (2,0) -- cycle;

        \fill[blue!25]
          (-2,0) -- (0,{2*\t}) -- (2,0) -- cycle;



        %------------------
        % 境界（点線）
        %------------------

        % 折れ線
        \draw[very thick,dashed] (-2,0)--(0,{2*\u})--(2,0);
        % 折れ線
        \draw[very thick,dashed] (-2,0)--(0,{2*\t})--(2,0);

        % x軸（|x|>=2）
        \draw[very thick,dashed] (-\xmax,0)--(-2,0);
        \draw[very thick,dashed] (2,0)--(\xmax,0);

        %------------------
        % 軸
        %------------------

        \draw[->,>=stealth,semithick] (-\xmax,0)--(\xmax,0) node[right]{\(x\)};
        \draw[->,>=stealth,semithick] (0,\ymin)--(0,\ymax) node[right]{\(y\)};

        % ラベル
        \node at (2.4,0.8) {\(U_t\)};
      \end{tikzpicture}
      \begin{tikzpicture}[scale=0.50]
        \def\t{0.7}
        \def\xmax{3.5}
        \def\ymin{-1}
        \def\ymax{3}

        %------------------
        % 境界（点線）
        %------------------

        % 折れ線
        \draw[ultra thick] (-2,0)--(0,{2*\t})--(2,0);
        %------------------
        % 軸
        %------------------

        \draw[->,>=stealth,semithick] (-\xmax,0)--(\xmax,0) node[right]{\(x\)};
        \draw[->,>=stealth,semithick] (0,\ymin)--(0,\ymax) node[right]{\(y\)};

        % ラベル
        \node at (1.4,1.2) {\(Z_t\)};
      \end{tikzpicture}
      %\caption{example \eqref{item:not-gentle}}
      %\label{fig:not-gentle}
    \end{figure}
\end{frame}

\begin{frame}
  \frametitle{うまくいかない例（続）}
    \(X\)の閉集合\(Z\)を次で定める．
    \[
      Z\coloneqq \left\{(x,y)\in X;\ (x+2)^{2}+(y-1)^{2}=1\right\}\setminus \left\{(-1,0)\right\}. 
    \]
    \(Z\)上の定数層の\(X\)への零による延長\(\kk_{Z}\)を考える．

    \(j\leqq0\)のとき，
    任意の実数\(s\leqq t\)と任意の点\(x\in Z_{s}\setminus U_{t}\)に対し\(
      H^{j}_{X\setminus U_{t}}(\kk_{Z})_{x}=0
    \). 

    しかし\(
      H^{0}(U_{1};\kk_{Z})=\kk
      \to 
      0 = H^{0}(U_{0};\kk_{Z})
    \)は単射ではない．
      \begin{figure}[htb]
      \centering
      \begin{tikzpicture}[scale=0.50]
        \def\t{0.9}
        \def\xmax{3.5}
        \def\ymin{-1}
        \def\ymax{3}

        %------------------
        % 塗りつぶし
        %------------------

        % y<0 の部分
        \fill[blue!25] (-\xmax,0) rectangle (\xmax,\ymin);

        % 三角形部分 y < t(2-|x|)
        \fill[blue!25]
          (-2,0) -- (0,{2*\t}) -- (2,0) -- cycle;

        %------------------
        % 境界（点線）
        %------------------
        % 折れ線
        \draw[ultra thick] (-2,0)--(0,{2*\t})--(2,0);
        
        % 折れ線
        \draw[very thick,dashed] (-2,0)--(0,{2*\t})--(2,0);

        % x軸（|x|>=2）
        \draw[very thick,dashed] (-\xmax,0)--(-2,0);
        \draw[very thick,dashed] (2,0)--(\xmax,0);

        %------------------
        % 軸
        %------------------

        \draw[->,>=stealth,semithick] (-\xmax,0)--(\xmax,0) node[right]{\(x\)};
        \draw[->,>=stealth,semithick] (0,\ymin)--(0,\ymax) node[right]{\(y\)};

        % 円
        \def\s{1} % 半径 t > 0
        \draw[thick] (-2,1) circle (\s);
        \fill[black!0](-2,0) circle (0.1);
        \draw (-2,0) circle (0.1);

        % ラベル
        \node at (2.2,0.8) {\(Z_t\)};
        \node at (2.4,-0.6) {\(U_t\)};
        \node at (-1.4,2.4) {\(Z\)};
      \end{tikzpicture}
      %\caption{example \eqref{item:not-gentle}}
      %\label{fig:not-gentle}
    \end{figure}
    \begin{alertblock}{工夫}
      「端点が移動する」という条件を付け加える．
    \end{alertblock} 
\end{frame}

\begin{frame}
    \frametitle{主結果2}
  \begin{THM}[大柴, 打切り版非特性変形補題]
    \(X\)をハウスドルフ空間とする．
    \(F\in\Dompl(\kk_X)\)とし，
    \(\left(U_{t}\right)_{t\in\rr}\)を\(X\)の開集合の族とする．
    \textcolor{brown}{\(k\in\zz\)とする．}
    \(
      Z_{s}\coloneqq \bigcap_{t>s}\overline{U_t\setminus U_s}\ (s\in\rr)
    \)とおく．
    次を仮定する．%\vspace{-8pt}
    \begin{enumerate}[(i)]
      \setlength{\leftskip}{10pt}
      \item \label{item:t-nonchar1}
      任意の\(t\in\rr\)に対し\(U_{t}=\bigcup_{s<t}U_{s}\)である．
      \item \label{item:t-nonchar2}
      任意の実数\(s\leqq t\)に対し\(
        \overline{U_{t}\setminus U_{s}}\cap \supp{F}
      \)はコンパクトである．
      \item \label{item:t-nonchar3}
      任意の実数\textcolor{brown}{\(s\in\rr\)}と
      任意の点\(x\in Z_{s}\setminus U_{\textcolor{brown}{s}}\)に対し，
      \textcolor{brown}{\(j\leqq k\)ならば}
      \(
        (H^{j}_{X\setminus U_{s}}(F))_{x}=0. 
      \)
      \textcolor{darkred}{\item \label{item:t-nonchar4} 
      任意の実数\(s< t\)に対し，\(Z_s\subset U_{t}\)が成り立つ．}
    \end{enumerate}
    このとき，任意の\(t\in\rr\)に対し，
    次の射は\textcolor{brown}{\(j<k\)ならば同型}であり，
    \textcolor{darkred}{\(j=k\)ならば単射}である．
    \[
      H^{j}\left(\bigcup_{s\in\rr}U_{s};F\right)
      \to
      H^{j}\left(U_{t};F\right).
    \]
  \end{THM}
\end{frame}


\begin{frame}
  \frametitle{仮定 (iv) をみたす例}
    \begin{center}
    \eqref{item:t-nonchar4}：任意の実数\(s< t\)に対し\(Z_s\subset U_{t}\)．
    \end{center}
    実数\(t\)に対し関数\(\varphi_{t}\)を次で定める．
    \(t\leqq 0\)のとき\(\varphi_{t}(x)\equiv 0\)とする．
    
    \(t> 0\)のときは以下のように定める．\[
      \varphi_{t}(x)\coloneqq\begin{cases}
        \sqrt{t^{2}-x^{2}} & (\left|x\right|<\textcolor{black}{t}),\\
        0&(\left|x\right|\geqq \textcolor{black}{t}).
      \end{cases}
    \]
    このとき，\(\rr^2\)の開集合\(U_{t}\)を 
    \(
      U_{t}\coloneqq
      \left\{(x,y)\in X;\ y<\varphi_{t}(x)\right\}
    \)で定めると\eqref{item:t-nonchar4}をみたす．
    \begin{figure}[htb]
      \centering
      \begin{tikzpicture}[scale=0.6]
        \def\t{1.5}
        \def\s{1}

        %------------------
        % 塗りつぶし領域
        %------------------

        % 半円の下側 y < sqrt(t^2 - x^2)
        \fill[green!25,domain=-\t:\t,samples=200]
          (-\t,0)
          -- plot(\x,{sqrt(\t*\t-\x*\x)})
          -- (\t,0)
          -- cycle;
        \draw[dashed, very thick](0,0) ellipse (1.5 and 1.5);

        % 半円の下側 y < sqrt(s^2 - x^2)
        \fill[blue!25,domain=-\s:\s,samples=200]
          (-\s,0)
          -- plot(\x,{sqrt(\s*\s-\x*\x)})
          -- (\s,0)
          -- cycle;
        \draw[dashed, very thick](0,0) ellipse (1 and 1);

        % y < 0 の部分
        \fill[blue!25] (-3.5,0) rectangle (3.5,-2);

        %------------------
        % 境界
        %------------------

        % 半円（上半分）
        %\draw[thick,domain=-\t:\t,samples=300]
        %  plot(\x,{sqrt(\t*\t-\x*\x)});

        % x軸（|x|>=t の境界として）
        \draw[dashed, very thick] (-3.5,0)--(-\s,0);
        \draw[dashed, very thick] (\s,0)--(3.5,0);

        %------------------
        % 軸
        %------------------

        \draw[->,>=stealth,semithick] (-3.5,0)--(3.5,0) node[right]{\(x\)};
        \draw[->,>=stealth,semithick] (0,-2)--(0,2) node[right]{\(y\)};

        % ラベル
        \node at (2.3,-0.6) {\(U_t\)};
      \end{tikzpicture}
      \begin{tikzpicture}[scale=0.6]
        \def\t{1}

        %------------------
        % 塗りつぶし領域
        %------------------

        % 半円の下側 y < sqrt(t^2 - x^2)
        \fill[black!0,domain=-\t:\t,samples=200]
          (-\t,0)
          -- plot(\x,{sqrt(\t*\t-\x*\x)})
          -- (\t,0)
          -- cycle;
        \draw[ultra thick](0,0) ellipse (1 and 1);

        % y < 0 の部分
        \fill[black!0] (-3.5,0) rectangle (3.5,-2);

        %------------------
        % 境界
        %------------------

        % 半円（上半分）
        %\draw[thick,domain=-\t:\t,samples=300]
        %  plot(\x,{sqrt(\t*\t-\x*\x)});

        % x軸（|x|>=t の境界として）
        %\draw[dashed, very thick] (-3.5,0)--(-\t,0);
        %\draw[dashed, very thick] (\t,0)--(3.5,0);

        %------------------
        % 軸
        %------------------

        \draw[->,>=stealth,semithick] (-3.5,0)--(3.5,0) node[right]{\(x\)};
        \draw[->,>=stealth,semithick] (0,-2)--(0,2) node[right]{\(y\)};

        % ラベル
        \node at (1.6,0.6) {\(Z_t\)};
      \end{tikzpicture}
      %\caption{example \eqref{item:gentle2}}
      %\label{fig:gentle2}
    \end{figure}
\end{frame}


\begin{frame}
    \frametitle{主結果2（再掲）}
  \begin{THM}[大柴, 打切り版非特性変形補題]
    \(X\)をハウスドルフ空間とする．
    \(F\in\Dompl(\kk_X)\)とし，
    \(\left(U_{t}\right)_{t\in\rr}\)を\(X\)の開集合の族とする．
    \textcolor{brown}{\(k\in\zz\)とする．}
    \(
      Z_{s}\coloneqq \bigcap_{t>s}\overline{U_t\setminus U_s}\ (s\in\rr)
    \)とおく．
    次を仮定する．%\vspace{-8pt}
    \begin{enumerate}[(i)]
      \setlength{\leftskip}{10pt}
      \item \label{item:t-nonchar1}
      任意の\(t\in\rr\)に対し\(U_{t}=\bigcup_{s<t}U_{s}\)である．
      \item \label{item:t-nonchar2}
      任意の実数\(s\leqq t\)に対し\(
        \overline{U_{t}\setminus U_{s}}\cap \supp{F}
      \)はコンパクトである．
      \item \label{item:t-nonchar3}
      任意の実数\textcolor{brown}{\(s\in\rr\)}と
      任意の点\(x\in Z_{s}\setminus U_{\textcolor{brown}{s}}\)に対し，
      \textcolor{brown}{\(j\leqq k\)ならば}
      \(
        (H^{j}_{X\setminus U_{s}}(F))_{x}=0. 
      \)
      \textcolor{darkred}{\item \label{item:t-nonchar4} 
      任意の実数\(s< t\)に対し，\(Z_s\subset U_{t}\)が成り立つ．}
    \end{enumerate}
    このとき，任意の\(t\in\rr\)に対し，
    次の射は\textcolor{brown}{\(j<k\)ならば同型}であり，
    \textcolor{darkred}{\(j=k\)ならば単射}である．
    \[
      H^{j}\left(\bigcup_{s\in\rr}U_{s};F\right)
      \to
      H^{j}\left(U_{t};F\right).
    \]
  \end{THM}
\end{frame}

\begin{frame}
    \frametitle{まとめ}
    \[\vcenter{\xymatrix@C=36pt@R=10pt{
        \txt{\(\muS[](F)\)\\
        \cite{KS82}}
        %\ar[dd]
        &
        \txt{超局所モースの補題\\
        \cite{KS90}\\\(H^{j}(M_b;F)\cong H^{j}(M_a;F)\)}
        %\ar[dd]
        &
        \txt{非特性変形補題\\
        \cite{KS90}\\
        \(H^{j}\left(\bigcup\limits_{s\in\rr}U_{s};F\right)
        \cong
        H^{j}\left(U_{t};F\right)\)}
        %\ar[dd]
        \\
        %\ar[r] 
        &
        &
        %\ar[l]
        \\
        \txt{\(\muS[k](F)\)\\
        \cite{KFS03a}}
        &
        \txt{打切り超局所モースの補題\\
        主結果1\\
        \(H^{j}(M_b;F)\to H^{j}(M_a;F)\)\\
        \(j<k\)で同型，\(j=k\)で単射}
        &\txt{打切り非特性変形補題\\主結果2\\\(
        H^{j}\left(\bigcup\limits_{s\in\rr}U_{s};F\right)
        \to
        H^{j}\left(U_{t};F\right)\)\\\(
        j<k\)で同型，\(j=k\)で単射}
        }}\]
    \vspace{10pt}
    \begin{center}
      結果：層のコホモロジーの延長可能性を次数に関して精密化した．
    \end{center}

\end{frame}

\appendix
\section*{参考文献}

\begin{frame}[allowframebreaks]{参考文献}
  \begin{thebibliography}{90}\beamertemplatetextbibitems
  \bibitem[FM07]{FM07} Teresa Monteiro Fernandes, Ana Rita Martins. 
  \textit{Truncated Microsupport and Hyperbolic Inequalities}, 
  Nagoya Mathematical Journal, 2007, Vol.\ 185, pp.63--91. 
  \bibitem[GKS12]{GKS12} St\'ephane Guillermou, 
  Masaki Kashiwara, Pierre Schapira, 
  \textit{Sheaf Quantization of Hamiltonian Isotopies and Applications to Nondisplaceability Problems}, 
  Duke. Math. J. 161, No. 2, pp.201--245, 2012.  
  \bibitem[Hor90]{Hor90} Lars H\"ormander, 
  \textit{The Analysis of Linear Pertial Differential Operators I}, Second Edition,
  Grundlehren der Mathematischen Wissenschaften, 256, Springer, 1989.
  \bibitem[KFS03a]{KFS03a} Masaki Kashiwara, Teresa Monteiro Fernandes, Pierre Schapira, 
  \textit{Truncated Microsupport and Holomorphic Solutions 
  of \(D\)-modules}, 
  Ann. Scient. \'Ec. Norm. Sup., s\'erie 4, tome 36, 
  pp.583--399, 2003.
  \bibitem[KFS03b]{KFS03b} Masaki Kashiwara, Teresa Monteiro Fernandes, Pierre Schapira, 
  \textit{Involutivity of Truncated microsupports}, 
  Bull. Soc. math. France, 131 (2), 2003, 
  pp.259--266.
  \bibitem[KS82]{KS82} Masaki Kashiwara and Pierre Schapira,
  \textit{Micro-support des faisceaux: applications 
  aux modules diff{\'e}rentiels,}
  C.~R.~Acad.\ Sci.\ Paris s{\'e}rie I Math 295 8, 487--490 France (1982).
  \bibitem[KS90]{KS90} Masaki Kashiwara, Pierre Schapira, 
  \textit{Sheaves on Manifolds}, 
  Grundlehren der Mathematischen Wissenschaften, 292, Springer, 1990.
  \bibitem[Sato69]{Sato69} Mikio Sato, 
  \textit{Hyperfuncitons and 
  Partial Differential equations}, 
  Proc. Int. Conf. on Functional Analysis, Tokyo, 1969.
  \bibitem[Sato70]{Sato70} Mikio Sato, 
  \textit{Regularity of Hyperfunciton Solutions of 
  Partial Differential equations}, 
  Actes, Congres intern. Math., 
  Tome 2, pp.785--794, 1970.
\end{thebibliography}
\end{frame}

\begin{frame}
  \frametitle{単射性に関する補足}
  \begin{THM}[Holmgrenの一意性定理 (例えば{\cite[Theorem 8.6.5]{Hor90}})]
    \(u\)を\(X\)上で（実解析係数）線形偏微分方程式\(Pu=0\)をみたす超関数とする．
    \(P\)に関する非特性\(C^1\)曲面の近傍を考える．
    曲面の片側において\(u=0\)ならば，近傍全体で\(u=0\)である．
  \end{THM}
  すなわち，初期面が\(P\)に関して非特性的なら，
  \(Pu=0\)の局所解は，あったとすればただ一つ！
  
  層係数コホモロジーの言葉で言うと，
  超関数解の層\(F=\RHOM_{\cD_{X}}(\cM,\cB_{M})\)に対し，
  \(x\in X\)の近くで
  \[
    H^{0}(B;F)\to H^{0}(B\cap \left\{\varphi<0\right\};F)
  \]
  が単射ということ．（\(B\)は\(x\)の開近傍，\(\varphi\)は近傍上の\(C^{1}\)関数）

  打切り超局所台の言葉で言うと，
  \[
    \muS[0](F)\subset \Char(\cM)\cap \cotb[M].
  \]
  ということ（主結果2から\textcolor{red}{一意性を半大域化}できる）．
\end{frame}

\begin{frame}
  \frametitle{仮定 (iv) をみたす例（追加）}
    \begin{center}
    \eqref{item:t-nonchar4}：任意の実数\(s< t\)に対し\(Z_s\subset U_{t}\)．
    \end{center}
    \(\rr^2\)の開集合\(U_{t}\)を次のように定める．
    \begin{itemize}
      \item \(t\leqq 0\)のときは\(U_{t}\coloneqq \left\{(x,y)\in X;\ y<0\right\}\). 
      \item \(t> 0\)のときは\(
        U_{t}\coloneqq
        \left\{(x,y)\in X;\ 
        \begin{array}{cc}
          \text{\(\left|x\right|<\tfrac{1}{t}\)ならば，}&y<t.\vspace{2pt}\\
          \text{そうでないとき，}&y<\tfrac{1}{\left|x\right|}.
        \end{array}
        \right\}\).
    \end{itemize}
    このとき，\((U_{t})\)は\eqref{item:t-nonchar4}をみたす．
    \begin{figure}[htb]
      \centering
      \begin{tikzpicture}[scale=0.6]
        \def\t{2}
        \def\s{1.5}
        \def\a{1/\t}
        \def\b{1/\s}
        % --- 塗りつぶし部分1 ---

        % 中央部分 |x|<1/t, y<t
        \fill[green!20]
          (-\a,0) rectangle (\a,\t);

        % 右側 |x|>=1/t
        \fill[green!20,domain=\a:3,samples=200]
          (\a,0)
          -- plot(\x,{1/\x})
          -- (3,0)
          -- cycle;

        % 左側 |x|>=1/t
        \fill[green!20,domain=-3:-\a,samples=200]
          (-3,0)
          -- plot(\x,{1/abs(\x)})
          -- (-\a,0)
          -- cycle;

        % --- 塗りつぶし部分2 ---

        % 中央部分 |x|<1/s, y<s
        \fill[blue!30]
          (-\b,0) rectangle (\b,\s);

        % 右側 |x|>=1/t
        \fill[blue!30,domain=\b:3,samples=200]
          (\b,0)
          -- plot(\x,{1/\x})
          -- (3,0)
          -- cycle;

        % 左側 |x|>=1/t
        \fill[blue!30,domain=-3:-\b,samples=200]
          (-3,0)
          -- plot(\x,{1/abs(\x)})
          -- (-\b,0)
          -- cycle;


        % --- 境界線 ---

        % 水平線 y=t
        \draw[dashed, thick] (-\a,\t)--(\a,\t);
        % 水平線 y=s
        \draw[dashed, thick] (-\b,\s)--(\b,\s);

        % 双曲線
        \draw[dashed, thick,domain=\a:3,samples=200]
          plot(\x,{1/\x});
        \draw[dashed, thick,domain=-3:-\a,samples=200]
          plot(\x,{1/abs(\x)});

        % 軸
        \draw[->,>=stealth,semithick] (-3.5,0)--(3.5,0) node[right]{\(x\)};
        \draw[->,>=stealth,semithick] (0,-0.1)--(0,3) node[right]{\(y\)};

        % 補助線
        %\draw[dashed] (\a,0)--(\a,2);
        %\draw[dashed] (-\a,0)--(-\a,2);

        % ラベル
        \node at (2.1,1) {\(U_t\)};
      \end{tikzpicture}
      \begin{tikzpicture}[scale=0.6]
        \def\t{1.5}
        \def\a{1/\t}

        % 軸
        \draw[->,>=stealth,semithick] (-3.5,0)--(3.5,0) node[right]{\(x\)};
        \draw[->,>=stealth,semithick] (0,-0.1)--(0,3) node[right]{\(y\)};

        % --- 塗りつぶし部分 ---

        % --- 境界線 ---

        % 水平線 y=t
        \draw[ultra thick] (-\a,\t)--(\a,\t);
        % ラベル
        \node at (1.2,\t) {\(Z_t\)};
      \end{tikzpicture}
      %\caption{example \eqref{item:gentle3}}
      %\label{fig:gentle3}
    \end{figure}
\end{frame}
\begin{frame}
  \frametitle{主結果2の証明の概要 (1/2)}
  \(\dip 
    H^{j}\left(\bigcup_{s\in\rr}U_{s};F\right)
    \to
    H^{j}\left(U_{t};F\right)
  \)の\(j<k\)での同型と\(j=k\)での単射性を示したい．

  ある補題\cite[Proposition 1.12.6]{KS90}を用いると，
  次を示すことに帰着される．
  \begin{align}
    \tag*{\((a)^j_s\)}\label{eq:t-nonchar-a}
    \mu_{s}^{j}&\colon
    \indlim[t>s] H^{j}(U_t;F) \simar H^{j}(U_s;F),\\
    \tag*{\((a')^k_s\)}\label{eq:t-nonchar-a'}
    \mu_{s}^{k}&\colon
    \indlim[t>s] H^{k}(U_t;F) \hookrightarrow H^{k}(U_s;F),\\
    \tag*{\((b)^j_s\)}\label{eq:t-nonchar-b}
    \lambda_{s}^{j}&\colon
    H^{j}(U_s;F)\simarr \prolim[r<s] H^{j}(U_r;F).
  \end{align}
  \ref{eq:t-nonchar-a}を示すには，次の射の列を考える．
  \begin{equation*}%\label{eq:t-nocha-d.t.}
    \RG(Z_{s}\cup U_{s};\mres[\RG_{X- U_{s}}(F)]{Z_{s}\cup U_{s}})
    \to
    \RG(Z_{s}\cup U_{s};\mres[F]{Z_{s}\cup U_{s}})
    \to
    \RG(U_{s};F)
    \overset{+1}{\longrightarrow}. 
  \end{equation*}
  仮定\eqref{item:t-nonchar3}から\(j\leqq k\)で第1項が\(0\)であることがわかる．
\end{frame}

\begin{frame}
  \frametitle{主結果2の証明の概要 (2/2)}
  したがって，次の制限射は\(j<k\)なら同型，\(j=k\)なら単射である．
  \[
    H^{j}(Z_{s}\cup U_{s};\mres[F]{Z_{s}\cup U_{s}})
    \to
    H^{j}(U_{s};F).
  \]
  次の補題を示せば，
  \ref{eq:t-nonchar-a}と\ref{eq:t-nonchar-a'}がわかる．
  \begin{LMM}%\label{clm:ind-nonchar}
    任意の\(j\in\zz\)と\(s\in\rr\)に対し，次の自然な射は同型である．
    \[
      \indlim[t>s]H^{j}(U_t;F)
      \cong 
      H^{j}(Z_{s}\cup U_{s};\mres[F]{Z_{s}\cup U_{s}}).
    \]
  \end{LMM}
  \ref{eq:t-nonchar-b}は\(j\)に関する帰納法とMittag-Leffler条件を用いることで示される．
\end{frame}
\begin{frame}
  \frametitle{証明の補足}

  ある補題\cite[Proposition 1.12.6]{KS90}というのは，
  柏原の補題と呼ばれる次の主張．
  \begin{PRP}[{\cite[Proposition 1.12.6]{KS90}}]
    \(
      \left(X_s,\rho_{s,t}\colon X_t\to X_s\right)_{(s,(s\leqq t))}
    \)を\(\rr\)で添字づけられた射影系とする．
    各\(s\in\rr\)に対し，次の自然な写像がどちらも単射（全射）であると仮定する．
    \[
      \lambda_s\colon X_s\to\prolim[r<s]X_r
      ,\quad
      \mu_s\colon \indlim[t>s]X_t\to X_s.
    \]
    このとき，すべての射\(
      \rho_{s_{0},s_{1}}\) \((s_0\leqq s_1)
    \)は単射（全射）となる．
  \end{PRP}
  これを用いると，
  \ref{eq:t-nonchar-a}と\ref{eq:t-nonchar-b}より，
  制限射
  \[
    H^{j}\left(U_{t};F\right)
    \to
    H^{j}\left(U_{s};F\right)\quad  (s\leqq t)
  \]
  が\(j<k\)で同型，\(j=k\)で単射となることがわかる．

  最後に，もういちどMittag-Leffler条件を用いることで証明が終わる．
\end{frame}

\begin{frame}
  \frametitle{展望：相対版打切り超局所モースの補題}
  \begin{itemize}
    \item \(f\colon Y\to X\)：多様体の射．
    \item \(\varphi\colon Y\to \rr\)：\(C^{1}\)級関数で
    \(Y_{t}\coloneqq \varphi^{-1}((-\infty,t))\), 
    \(\widetilde{Y}_{t}\coloneqq \varphi^{-1}((-\infty,t])\)とし，
    \(\iota_{t}\colon Y_{t}\hookrightarrow Y\).
    \item \(G\in\Dompb(Y)\)で
    \(f_{t}\colon \supp(G)\cap \overline{Y_t}\to X\)はすべての\(t\in\rr\)に対し固有．
  \end{itemize}
  \(k\in\zz\)とし\(t_{\circ}\in\rr\)とする．
  \(y\in Y\setminus \widetilde{Y}_{t_{\circ}}\)ならば
  次が成り立つとする．\[
      d\varphi(y)\notin \left(
        \muS[k](G)+ \mdiff[f]\left(
          Y\pb[X]\cotb[X]
        \right)
      \right).
    \]
  このとき，\(t>t_{\circ}\)に対して，次が成り立つと期待される．
  \begin{CNJ}
    次の層の射は\(j<k\)ならば同型，\(j=k\)ならば単射である．
    \[
      \Rder[j]{f_{\ast}}G
      \to
      \Rder[j]{{f_{t}}_{\ast}}\iota^{-1}_{t}G.
    \]
    ただし\(f_{t}=f\circ \iota_{t}\)とおいた．
  \end{CNJ}
\end{frame}

\begin{frame}
  \frametitle{主結果2の設定について}
  修論中では「任意の開集合がパラコンパクト」を課していたが，不要であることが分かった．

  \begin{alertblock}{注意}
    \(X\)を位相空間とする．
    一般の部分空間\(S\)と自然な連続写像\(
      i\colon S\hookrightarrow X
    \)を考える．
    \(F\in\Dompl(\kk_{X})\)とする．
    自然な射\[
      \Rder[]\left(\Gamma(S;\blk)\circ i^{-1}\right)(F)
      \to
      \Rder[]{(\Gamma(S;\blk))}\circ i^{-1}(F)
    \]は一般には同型でない．    
  \end{alertblock}
  次が言えることがわかったので，今の設定では問題ない．
  \begin{PRP}
    \(X\)をハウスドルフ空間とする．
    \(U\subset X\)を開集合，\(Z\subset X\)をコンパクト集合とする．
    \(U\cap Z=\varnothing\)ならば，
    \(F\in\Dompl(\kk_{X})\)と\(S=U\cup Z\)に対し上の射は同型．
  \end{PRP}
\end{frame}

\begin{frame}
  \frametitle{主結果2の設定について（続）}
  \begin{PRP}
    \(X\)をハウスドルフ空間とする．
    \(U\subset X\)を開集合，\(Z\subset X\)をコンパクト集合とする．
    \(U\cap Z=\varnothing\)ならば，
    \(F\in\Dompl(\kk_{X})\)に対し次の射は同型．\[
      \Rder[]\left(\Gamma(U\cup Z;\blk)\circ i^{-1}\right)(F)
      \to
      \Rder[]{(\Gamma(U\cup Z;\blk))}\circ i^{-1}(F).
    \]
  \end{PRP}
  これの証明の方針は，「少し膨らませてからMayer--Vietris三角を比較」．
  すなわち，
  \[
    \Rder(\Gamma(U\cup Z;\blk)\circ i^{-1})(F)
    \to \RG(U;F)\oplus\RG(Z;F)
    \to \RG(Z;\RG_{U}F)
    \overset{+1}{\longrightarrow}
  \]
  と\(\mres[F]{U\cup Z}\)に対して定まる三角
  \[
    \RG(U\cup Z;\mres[F]{U\cup Z})
    \to \RG(U;F)\oplus\RG(Z;F)
    \to \RG(Z;\RG_{U}F)
    \overset{+1}{\longrightarrow}
  \]を比較することで\(
    \Rder(\Gamma(U\cup Z;\blk)\circ i^{-1})(F)
    \simar
    \RG(U\cup Z;\mres[F]{U\cup Z})
  \)となる．
\end{frame}


\begin{comment}
\begin{frame}[noframenumbering]
aaa
\end{frame}
\end{comment}

\end{document}